\documentclass[11pt,a4paper]{article}
\usepackage[margin=2cm, top=1.5cm]{geometry}
\usepackage{times}
\usepackage{graphicx}
\usepackage[none]{hyphenat}
\setlength{\parindent}{0pt}
\setlength{\parskip}{0.65em}

\begin{document}

\begin{center}
\Large \textbf{Statement of Purpose}\\
\vspace{0.2cm}
\normalsize Arman Asgharpoor\\
RAINZ CDT Application -- University of Oxford\\
\vspace{0.1cm}
\textbf{Word Count: 959}
\end{center}

\vspace{0.1cm}

Dear Admission Committee,

Five minutes of sensor data from our lab produced the reconstruction
in Figure~\ref{fig:mesh}. The geometry is there, but the map has no
memory and no meaning. It cannot tell what any object is, whether it
was there last week, or whether something small has changed since the
last scan. Building a representation that accumulates spatial knowledge
over time and resolves changes at the scale of individual objects from
camera data alone is the problem I want to work on.

Last November, a Unitree Go2 navigated our lab using this map. Localization held to fifteen centimeters, but when a person stood near the glass partition, the robot stopped. It could not separate the person from the wall. Both entered the map as permanent obstacles. The problem was not the sensor or the algorithm. The map had no concept of object identity or temporal persistence. Robots inspecting wind turbines, solar arrays, or nuclear facilities face this gap at higher stakes. Prior work such as VILENS~\cite{wisth2022vilens} provides robust, low-drift multi-sensor odometry (visual–inertial–lidar–leg) that enables accurate high-frequency local terrain/elevation mapping for planning and control, and LEXIS~\cite{kassab2024lexis} shows semantics can be embedded into such maps through CLIP features and label propagation. However, these approaches primarily enrich map-level geometry (e.g., terrain/surface representations) rather than producing object- or instance-centric semantic maps. Neither can tell whether a specific object has changed, disappeared, or degraded. A facility valve inspected months apart will differ geometrically from lighting, oxidation, and surface accumulation, yet it is the same object in a changed state. That is the gap I want to close.

\begin{figure}[h]
\centering
\includegraphics[width=0.75\textwidth]{mesh.png}
\caption{Dense 3D reconstruction of our lab.}
\label{fig:mesh}
\end{figure}

My Master's research developed a sampling-rate-aware neural network
for IMU attitude estimation: a 1D-CNN feature extractor with a
bidirectional LSTM and an explicit sampling-rate fusion layer, trained
with Gaussian noise injection across seven public datasets covering
200$+$ km of motion. It reduced orientation error by 37\% without
retraining across sensors, published in
\textit{Measurement}~\cite{golroudbari2023imu}. At Fasta Robotics, I
integrated ORB-SLAM3 with ROS2 on RealSense D435i and ZED 2i across multi-floor warehouse layouts, trained YOLOv12 for human-aware navigation reaching 95\% recall at
30\,FPS, and tracked people with DeepSORT. When someone moved behind a
forklift, the tracker lost them because SLAM held no semantic state to
help. The failure was not in either component. It was in the absence
of a shared representation between them.

With only 204 labeled images across 53 classes, I used CLIP embeddings
to keep object identity stable across viewpoint and lighting shifts,
reaching 81\% mAP. In TE-PINN, I embedded rotational kinematics on
SO(3) directly into the transformer training objective, so geometric
consistency was enforced by construction rather than hoped for through
the loss. The next step is reasoning about 3D object identity and change
from camera observations across time. 

Feed-forward geometric reconstruction backbones such as 
VGGT~\cite{wang2025vggt} and MapAnything~\cite{keetha2025mapanything} 
directly predict multi-view 3D structure in a single pass, with 
MapAnything additionally targeting globally consistent metric 
reconstructions from flexible and partial inputs. I intend to build 
on these models as geometric priors and integrate them with 
multi-session visual SLAM to reason about object identity and change 
across revisits separated by weeks. For that to be deployable over 
months, the map cannot grow without bound: PlanarMesh~\cite{wang2025planarmesh} 
shows that adaptive-resolution meshing can substantially reduce map 
footprint while retaining the detail needed to detect small changes.


Those models provide the geometric foundations, but not a principled 
way to reason about what has changed and what has not. A map built 
from one session carries geometry but no model of what it does not 
yet know. How a system should reconcile conflicting observations 
across sessions and update its belief is a question in probabilistic 
inference, not engineering. I have not had formal training to address 
that from first principles. I placed first among 120+ participants in the Oxford Machine Learning 
Summer School competition, received an Outstanding Reviewer recognition 
from IEEE, and have reviewed for ICRA and ICLR, so the research 
exposure is there. The grounding 
in estimation theory is not, and that is why a PhD is the right step 
now rather than more deployment. The DRS group is where that formal 
training connects directly to deployment on real platforms, which is 
the combination I cannot find elsewhere.


After the PhD I want to work on inspection systems for energy
infrastructure that are reliable enough to operate without human
verification of every output. RAINZ is structured around that goal. The cohort model means working on a shared industrial challenge across wind, solar, and nuclear rather than isolated academic problems. When a method fails in one energy context, the cohort sees whether the failure generalizes or is domain-specific. That cross-sector exposure is how you learn what makes a system deployable versus what makes it publishable. The industrial engagement and commercialisation training mean the gap between research and field deployment stays visible throughout. I want to build tools operators actually use, and I want to understand deeply enough why they fail to fix them when they do.

I am applying to work with Professor Maurice Fallon on the Long-Term 
3D Change Detection of Industrial Facilities project. My work has 
followed the same direction independently: from inertial state 
estimation, into visual-inertial SLAM, and toward the point where 
geometric representations fail because they carry no model of object 
identity or change over time. That is where this project begins. 
The Year 1 MSc at Manchester addresses the formal gap I identified 
through deployment rather than reading. Years 2 to 4 at Oxford, 
within a group that has already demonstrated these foundations in 
real facilities, is the environment where that gap can be properly 
closed. I want to build tools operators actually use, in infrastructure 
where the cost of a missed change is real.


\bibliographystyle{ieeetr}
\bibliography{ref}

\end{document}